\chapter*{Abstract} 
\addcontentsline{toc}{chapter}{Abstract}
\pagestyle{plain}
 
Steganalysis has typically relied on hand-crafted statistical means of discovering concealed data; 
however, in recent years the effort has largely shifted towards automated feature extraction techniques 
via deep learning. 
Since criminal and other malicious actors have been found to use steganographic methods in their line of work, 
reliable detection of such hidden content is of direct relevance to law enforcement 
and cybersecurity practitioners tasked with uncovering covert communication channels in digital forensic 
investigations. 
This thesis examines whether the newer Convolutional Neural Networks (CNNs) are more 
effective than traditional RS-analysis methods to detect Least Significant Bit (LSB) steganography 
utilized to hide data in lossless PNG image files. 
To evaluate these two detection methods, a quantitative experiment was performed using the BOSSBase 1.01 
dataset on 4 different embedding techniques: Sequential LSB, Randomized LSB, Edge-Adapted LSB, 
and LSB Matching, at three different 
payload capacities of 25\%, 50\%, and 75\%. 
The results of this investigation show that traditional RS-analysis 
or modern CNN methods can be optimal based on what embedding strategy has been used. 
RS-analysis was superior to CNN (SRNet) for both Sequential and Randomized LSB modifications across 
all examined payloads, achieving greater than 95\% accuracy for both modified versions even at the 
lowest embedding rate of 25\%. 
Conversely, for low payload Edge-Adaptive LSB, CNN demonstrated superior robustness with an 
accuracy of 79.12\% versus the RS-analysis accuracy of 72.50\% at a 25\% payload. 
Finally, only the CNN was capable of detecting LSB Matching (with an accuracy of 91.36\%), 
where the RS-analysis performed below chance level in detecting this embedding method. 
Therefore, these results demonstrate that traditional statistical techniques continue to be 
effective and sufficient for the basic Least Significant Bit substitution; however, deep 
learning techniques are needed for detecting Edge-Adaptive LSB and/or LSB Matching steganography.
